\section{数据说明与变量设计}
\subsection{数据来源与口径统一}
本文数据主要来自CEADs省级表观碳排放清单与国家统计局年度数据库，并辅以生态环境经济核算报告和公开统计资料进行口径核对。由于碳排放原始数据按省份发布，本文先按年份进行全国汇总，构建国家级碳排放总量序列；宏观解释变量则直接采用国家级年度指标，以保证时间序列的一致性和可比性。

数据源于 CEADs 省级清单与国家统计局数据库，辅以生态环境经济核算报告进行核对。因变量$C_t$来源于 CEADs 碳排放清单。该清单包含化石燃料燃烧和过程排放等信息，本文在同一年份内对省级排放值求和，得到国家总量。解释变量包括人口、GDP、人均GDP、城镇化率、第二产业占比、能源消费总量、煤炭消费占比、创新与治理代理指标等。通过该组合，本文从人口、经济、结构、能源、技术和治理六个维度描述碳排放驱动机制。

在国家级研究框架下，口径统一比变量数量更重要。为此本文遵循三项原则：第一，同一变量仅保留一个统计口径，避免不同来源混用；第二，优先采用连续年份完整、可追溯的数据源；第三，必要时对单位进行统一换算并在模型前进行一致性检查。该处理有助于降低“口径漂移”造成的估计偏差。

\subsection{变量构造与指标定义}
设$t$为年份，$C_t$为全国碳排放总量，$P_t$为总人口，$A_t$为富裕度，核心变量定义如下：
\begin{align}
C_t &= \sum_{i=1}^{N} C_{it}, \\
A_t &= \frac{GDP_t}{P_t}, \\
EI_t &= \frac{EC_t}{GDP_t}, \\
CI_t &= \frac{C_t}{GDP_t}, \\
PC_t &= \frac{C_t}{P_t}.
\end{align}

其中，$EC_t$为能源消费总量，$EI_t$表示能源强度，$CI_t$为碳强度，$PC_t$为人均碳排放。能源结构变量以煤炭消费占比表示，记为$CS_t$；城镇化变量以城镇人口占比表示，记为$UR_t$；产业结构变量以第二产业增加值占比表示，记为$IS_t$；技术进步变量以有效发明专利和科研机构规模的组合代理表示，记为$TP_t$；环境治理变量以污染治理投资强度表示。

本文在变量选取上注重理论与数据的匹配，所选指标均可从公开统计口径获取，且保证了跨年份的量纲一致，以增强实证结果的可靠性。相较于引入大量高维指标，本文更强调核心变量的稳定测度，以保证小样本时序估计的有效性。

\subsection{变量预处理与统计策略}
为了提高估计稳定性，本文对规模变量（如人口、富裕度、能源强度）采用对数变换，使系数可解释为弹性并缓解异方差。比例变量（如煤炭占比、城镇化率、产业结构占比）保留原始比例形式，以便直接解释结构变化的边际影响。对于数量级差异显著的变量，在进入机器学习模型前进行标准化处理，用于提升训练稳定性。

在统计分析流程上，本文先进行描述统计与相关性检验，再进入机制回归与预测比较，最后进行情景模拟。该顺序可确保每一步结论都建立在前一环节的可验证结果之上，避免直接从模型输出跳到政策结论。

\subsection{描述性统计特征}
描述性统计结果表明，在2005—2022年观测区间内，我国碳排放总量均值约30566.50 Mt，最小值17516.39 Mt，最大值40358.26 Mt，呈现显著长期上行特征。同期GDP均值约654617.21亿元，人均GDP均值约4.73万元，城镇化率均值约54.90\%，第二产业占比均值约42.24\%，能源消费总量均值约417734.00万吨标准煤，煤炭占比均值约64.72\%。

这些统计特征共同说明：一方面，经济扩张和能源总量增长仍是排放上行的重要背景；另一方面，产业和能源结构已出现持续调整，为后续达峰路径优化提供了可能。技术创新代理变量在样本期内显著增长，提示创新要素积累正在形成中长期减排潜力，但其效应释放存在时滞。

\subsection{缺失值处理与样本边界}
本文保持2005—2024年研究时段不变。根据数据可得性，碳排放观测值覆盖2005—2022年，2023—2024年仅在预测与情景模拟中用于趋势衔接，不纳入机制回归和描述统计的观测样本。也即，实证估计基于真实观测区间，预测值只用于外推比较与路径展示。

该处理的意义在于：既保证图表与情景分析的时间连续性，又避免将模型推断误作事实数据。为强化边界清晰性，本文在相关章节始终区分“观测值”“预测值”和“情景值”三类数据，并在解释结论时仅对观测样本作统计推断。

\subsection{数据质量控制与稳健性说明}
为降低测度误差风险，本文进行了跨来源口径核对、异常年份回查以及结果一致性核验三类质量控制，如遇数据冲突，则重新核对原始来源并调整模型设定。

需要指出，本文并不以“覆盖全部可能变量”为目标，而是以“核心机制可被稳定识别”为优先原则。这一取舍有助于在有限数据可得性条件下保证结论可靠性，并为后续研究在更长样本和更细粒度数据基础上的扩展留出空间。
